% !TEX root = ../../../main.tex
% 所属：第一部分：单变量分析学 / 第1章：实数理论——从Peano公理到实数的连续性
% 节：1.11 复数域
% 正文按原文件顺序保存；由 ../chapter1.tex 调用。
\section{复数域}
\begin{definition}
	
	在 $ \mathbb{R} \times \mathbb{R}$  上定义如下等价关系 
	$$(a, b) \sim(c, d):=a=c \text { 且 } b=d$$
	$$\mathbb{C}:=\mathbb{R} \times \mathbb{R} / \sim$$
	
	
	设 $x=(a, b), y=(c, d)$ 是两个复数．我们定义：
	
	$$
	\begin{aligned}
		x+y & =(a+c, b+d) \\
		x y & =(a c-b d, a d+b c)
	\end{aligned}
	$$
	$(\C,+,\cdot)$是一个域，$(1,0)$视作$1$,$(0,0)$视作$0$
	$$\dfrac{1}{x}:=(\dfrac{a}{a^2+b^2},-\dfrac{b}{a^2+b^2})$$
	规定$(0,1)=i,(0,1)\cdot(0,1)=i^2=(-1,0)$
	
	习惯上，记$(a,b)=a+bi$
\end{definition}
\begin{definition}
	如果 $a, b$ 是实数，而 $z=a+b i$ ，就把复数 $\bar{z}=a-b i$ 叫做 $z$ 的共轭数.数 $a, b$ 分别叫做 $z$ 的实部和虚部.
	
	有时写成
	
	$$
	a=\operatorname{Re}(z), \quad b=\operatorname{Im}(z)
	$$
	
	
\end{definition}
\begin{proposition}
	\label{proofcon}
	如果 $z$ 及 $w$ 是复数，那么
	
	（a）$\overline{z+w}=\bar{z}+\bar{w}$ ，
	
	（b）$\overline{z w}=\bar{z} \cdot \bar{w}$ ，
	
	（c）$z+\bar{z}=2 \operatorname{Re}(z), \quad z- \bar{z}=2 i \operatorname{Im}(z)$ ．
	
	（d）$z \bar{z}$ 是实数且是正数（除了 $z=0$ 时）
\end{proposition}

\begin{proof}
	证（a），（b），（c）都十分明显，今证（d）．将 $z$ 写成 $z=a+b i$ ，并注意 $z \bar{z}= a^2+b^2$ .就行了.
\end{proof}
\begin{definition}
	如果 $z$ 是一复数，它的绝对值 $|z|$ 是 $z \bar{z}$ 的非负平方根；即 $|z|=(z \bar{z})^{1 / 2}$ .
	
\end{definition}

\begin{remark}
	当 $x$ 是实数时 $\bar{x}=x$ ，因此，$|x|=\sqrt{x^2}$ ．于是，如果 $x \geqslant 0,|x|=x$ ：如果 $x<0,|x|=-x$
\end{remark} 
\begin{proposition}
		设 $z$ 和 $w$ 都是复数，那么
	
	（a）$|z|>0$ 除非 $z=0,|0|=0$ ，
	
	（b）$|\bar{z}|=|z|$ ，
	
	（c）$|z w|=|z||w|$ ，
	
	（d）$|\operatorname{Re} z| \leqslant|z|$ ，
	
	（e）$|z+w| \leqslant|z|+|w|$ ．
\end{proposition}

\begin{proof}
	（a）及（b）是显然的．设 $z=a+b i, w=c+d i, a, b, c, d$ 为实数．
	
	于是
	
	$$
	\begin{aligned}
		|z w|^2 & =(a c-b d)^2+(a d+b c)^2 \\
		& =\left(a^2+b^2\right)\left(c^2+d^2\right)=|z|^2|w|^2
	\end{aligned}
	$$
	
	
	或 $|z w|^2=(|z||w|)^2$ .由定理$1.8.12$所说的惟一性，便能推得（c）.
	
	今证（d），注意 $a^2 \leqslant a^2+b^2$ ，由此，
	
	$$
	|a|=\sqrt{a^2} \leqslant \sqrt{a^2+b^2} .
	$$
	
	
	今证（e），注意 $\bar{z} w$ 是 $z \bar{w}$ 的共轭数，所以 $z \bar{w}+\bar{z} w=2 \operatorname{Re}(z \bar{w})$ .由此，
	
	$$
	\begin{aligned}
		|z+w|^2 & =(z+w)(\bar{z}+\bar{w})=\bar{z}+z \bar{w}+\bar{z} w+w \bar{w} \\
		& =|z|^2+2 \operatorname{Re}(z \bar{w})+|w|^2 \\
		& \leqslant|z|^2+2|z \bar{w}|+|w|^2 \\
		& =|z|^2+2|z||w|+|w|^2=(|z|+|w|)^2 .
	\end{aligned}
	$$
	
	
	两端开平方即得（e）.
\end{proof}
\begin{proposition}
	(Cauchy-Schawtz不等式)
	如果 $a_1, a_2, \cdots, a_n$ 及 $b_1, b_2, \cdots, b_n$ 都是复数，那么
	
	$$
	\left|\sum_{j=1}^n a_j \bar{b}_j\right|^2 \leqslant \sum_{j=1}^n\left|a_j\right|^2 \sum_{j=1}^n\left|b_j\right|^2 .
	$$
	\label{Cauchy-Schawtz不等式}
\end{proposition}

\begin{proof}
	设 $A=\Sigma\left|a_j\right|^2, B=\Sigma\left|b_j\right|^2, C=\Sigma a_j \bar{b}_j$（在本证明的所有和中，$j$ 遍历 $1, \cdots, n$ 诸值）.
	
	如果 $B=0$ ，那么 $b_1=\cdots=b_n=0$ ，而定理的结论是明显的.
	
	因此，可假定 $B>0$ ．由定理 \ref{proofcon}，有
	
	$$
	\begin{aligned}
		\sum\left|B a_j-C b_j\right|^2 & =\sum\left(B a_j-C b_j\right)\left(B \bar{a}_j-\bar{Cb_j}\right) \\
		& =B^2 \sum\left|a_j\right|^2-B \bar{C} \sum a_j \bar{b}_j-B C \sum \bar{a}_j b_j+|C|^2 \sum\left|b_j\right|^2 \\
		& =B^2 A-B|C|^2 \\
		& =B\left(A B-|C|^2\right) .
	\end{aligned}
	$$
	
	
	因第一个和中的每一项非负，而知
	
	$$
	B\left(A B-|C|^2\right) \geqslant 0 .
	$$
	
	
	由 $B>0$ 推得 $A B-|C|^2 \geqslant 0$ ．这就是要证的不等式．
\end{proof}
